% ============================================================
%  H. Høffding, "Filosofien og Darwinismen"
%  Et Foredrag, holdt i Studenterforeningen den 21. Marts 1874
%  Nær og Fjern, nr. 93 (12. april 1874) og nr. 94
%  TRANSCRIPTION (Danish)
%  Source: md draft (hoffding-darwin.md); to be verified against
%  scan-1.pdf (Part I) and scan-2.pdf (Part II).
% ============================================================
\documentclass[12pt,a4paper]{article}
\usepackage[utf8]{inputenc}
\usepackage[T1]{fontenc}
\usepackage[danish]{babel}
\usepackage{microtype}
\usepackage{setspace}
\usepackage[margin=1.2in]{geometry}
\usepackage{fancyhdr}
\usepackage{hyperref}

\hypersetup{
  pdftitle={Filosofien og Darwinismen},
  pdfauthor={H. Høffding},
  hidelinks
}

\pagestyle{fancy}
\fancyhf{}
\rhead{Høffding, \emph{Filosofien og Darwinismen} (1874)}
\cfoot{\thepage}

\setstretch{1.3}

\title{\textbf{Filosofien og Darwinismen}\\[0.4em]
\large Et Foredrag, holdt i Studenterforeningen den 21.\ Marts}
\author{Dr.\ phil.\ Harald Høffding}
\date{\emph{Nær og Fjern}, nr.\ 93--94, 1874}

\begin{document}
\maketitle
\thispagestyle{fancy}

\bigskip

\noindent Det er en gammel Erfaring, at Modsætninger sætter Eftertanken
i Bevægelse. Hvad man før opfattede med Ligegyldighed og uden
Opmærksomhed, tildrager sig paa een Gang Interessen, naar man faaer
Øje paa de Modsætninger, der ere forbundne med det, og som det giver
Anledning til. Den, som ikke af andre Grunde vilde føle sig foranlediget
til at anke over Darwinismen, vilde faae et Motiv dertil ved de
modstridende Opfattelser af den, de diametralt modsatte Anskuelser om
dens Værd og Betydning, som man daglig kan støde paa. For den Ene er
Darwinismen en vild Spekulation, et Produkt af vor materialistiske Tid,
der i selvdannede Hypotheser søger Stadfæstelse paa sine underlige og
endelige Drifter og Tilbøjeligheder; for den Anden er den en genial og
storartet Anskuelse, der aabner store Perspektiver, tilfredsstiller
Menneskeaandens ideale Higen efter det Uendelige og vækker Begejstringen
for højere, fremskridende Udvikling. Og gaa vi tilbage til den snevrere
Kreds, hvor denne Lære er opstaaet, og hvor den videnskabelige Prøvelse
hører hjemme, til Naturforskerne selv, vil man ligeledes finde en
paafaldende Modsætning mellem Opfattelserne. Medens en stor Mængde
Naturforskere (især Englændere og Tydskere) anse Darwins Lære for i den
Grad sandsynlig, at de ere tilbøjelige til at overse det Manglende i
Bevisførelsen, ja at endog Flere indlade sig paa vidtgaaende Spekulationer
og forlade Fagvidenskabens faste Grund --- saa hævde Andre (navnlig
franske Forskere, til hvilke de fleste danske synes at slutte sig), at
en Forklaring ikke er et Bevis, men forudsætter et Bevis for at have
Gyldighed, og at det, trods al den Betydning, Hypotheser kunne have for
Videnskaben, dog er aldeles utilstedeligt noget Øjeblik at glemme
Hullerne i Bevisførelsen over det Tiltalende i den hypothetiske
Forklaring.

Der er altsaa god Grund til at anstille en Prøvelse af Darwinismens
logiske og ethiske Værd for at komme til en Afgjørelse af, hvilken af
disse stridende Domme der er den sande, og paa hvilken Maade vi skulle
stille os til denne mærkelige Anskuelse, som fra Videnskabsmandens
stille Kammer nu er trængt ud paa Gader og Stræder og anvendes som et
mægtigt Vaaben i vor Tids aandelige Kampe. Filosofien har, afseet fra
disse ydre Foranledninger, i Følge sin egen Opgave den Pligt at anstille
en saadan Prøvelse. Det er Filosofiens Opgave at blive sig den
videnskabelige Erkjendelses Væsen og Methode bevidst, at skærpe Indsigten
i de Stadier, en Anskuelse maa gjennemgaa, førend den kan indlemmes i de
videnskabelige Sandheders Rige, og at hævde de nødvendige Betingelser
for Erkjendelsens Sandhed og Gyldighed. Men dens Kald er ikke begrændset
hertil. Idet den undersøger det menneskelige Bevidsthedsliv, dets Love
og dets Elementer, og forfølger dets Udvikling gjennem Historien, søger
den tillige at indordne ethvert Erkjendelsesresultat som Led i Aandens
hele Organisme og maa derfor paa ethvert Punkt prøve, hvorledes det
vundne Resultat forholder sig som Bidrag til en almindelig
Verdensanskuelse. Intet Resultat er i Sandhed tilegnet af Menneskeaanden,
førend denne Prøvelse er anstillet. Og først naar dette er gjort, kan
Resultatet faae berettigede praktiske Følger, idet det tredie Led i den
filosofiske Undersøgelse bliver: hvilken praktisk og ethisk Betydning
den tilegnede Sandhed faaer.

Det er klart, at den logiske Prøvelse her er den vigtigste. Først maa
det dog afgjøres, om en Lære er sand eller ikke, førend man skrider til
at prøve dens Betydning i metafysisk og ethisk Henseende. Man er ofte
for tilbøjelig til at gaa den omvendte Vej og mene, at hvad der har
ideal Betydning, ogsaa maa have real Virkelighed. Af denne Tilbøjelighed
ere de theologiske og spekulative Lærdomme udsprungne, som beherske den
almindelige Bevidsthed. Og mange af Darwinismens Forkæmpere have stillet
sig paa samme Standpunkt; de have dannet en ny Theologi, et nyt Dogme af
den rent empiriske Hypothese, Darwin har opstillet. Erkjendelseslæren maa
opponere mod begge disse Retninger og søge først og fremmest at fæste
Blikket paa Sagens virkelige Sammenhæng. Den indskærper os Resignationens
Nødvendighed. Et gammelt Ordsprog siger: »Den, som troer, haster ikke.«
Man kunde med lige saa megen Ret sige: »Den, som veed, haster ikke!«
idet vi ved den Vidende forstaa den, der har Blik for Erkjendelsens
væsenlige Betingelser saa vel som for de Ideer, der føre Menneskeaanden
fremad under dens theoretiske og praktiske Stræben. Hvorfor skulde disse
Ideer nødvendigvis have en --- overnaturlig eller naturlig --- Virkelighed?
Ere de ikke netop Ideer ved at vise ud over det faktisk Givne, ved at
stille nye Opgaver for Tænken og Handlen? Den, hvis Verdensanskuelse
ikke fordrer en dogmatisk Fastslaaen af de højeste Ideers Indhold som
ydre Faktum, vil med større Ro og Upartiskhed lade den videnskabelige
Methode for Erkjendelsen af det Virkelige ske Fyldest og ganske bøje
sig for dens Resultater.

\section*{I.}

Man kan sige, at Videnskaben vilde staa stille uden Hypothesen. For at
finde den Lov, hvorefter en vis given Faktor virker, eller for at finde
den Aarsag, der ligger til Grund for en vis given Lovmæssighed i
Fænomenerne, maa man nødvendigvis vove sig ud over det umiddelbart
Foreliggende, maa ofte gjøre store Spring til helt andre Kredse af
Fænomener for at prøve, om den søgte Forklaring ikke findes der. Med
Rette har man sagt, at de Forskere, der staa stille og vente, indtil de
udenfra drives til at gaa videre, aldrig ville gjøre de store Opdagelser.
Opdagelsen forudsætter en Aktivitet, som gaaer ud over den blotte
Opsamlen af Fakta, som søger en Enhed i Mangfoldigheden, en fælles Lov
for alle Forskjellighederne. Det er ikke mindst denne Aktivitet og
Energi, der gjør Kopernikus, Kepler og Newton til saa store Skikkelser i
Videnskabens Historie. Men dette er kun den ene Side af Sagen. Den Vej,
Tanken gaaer fra Fænomenerne til Loven eller Aarsagen, er maaske den,
der tager sig dristigst og mest storartet ud; men det er dog ikke den,
der videnskabelig seet er den afgjørende. Vi kunne beundre hin Energi,
men vi stole ikke paa, hvad den udretter, naar den ikke formaaer at gaa
den modsatte Vej tilbage igjen, naar den ikke fra den antagne Lov, den
postulerede Aarsag forstaaer at føre os tilbage igjen til den fænomenale
Sammenhæng. Det vil med andre Ord sige: Hypothesen maa verificeres,
bekræftes ved sin Overensstemmelse med Kjendsgjerninger, der ere vundne
ved umiddelbar Iagttagelse. Ogsaa dette har Newton havt Blik for, naar
han fordrer, at den Aarsag, man antager til Forklaring af visse
Fænomener, maa være en \textit{vera causa}, o: en saadan, som i
Virkeligheden viser sig at have saadanne Virkninger, som man tillægger
den.

Hvad er nu det, Darwins Hypothese vil forklare? Det er den
Kjendsgjerning, at de levende Væsener sondre sig i visse bestemte Grupper
med særegne Karaktermærker, som for hver enkelt Gruppe staa fast og
vedvare --- altsaa de saakaldte Artsforskjelligheders Kjendsgjerning.
Plante- og Dyreverdenen sondre sig for den iagttagende og klassificerende
Forsken i skarpt afgrændsede Arter. Botanikerne og Zoologerne have efter
den almindelige Opfattelse endt deres Hverv, naar de ved omhyggelig
Undersøgelse af en Plantes eller et Dyrs indre og ydre Egenskaber have
indseet, hvilken Art, Slægt, Familie den maa henføres til. Man følger det
enkelte Individs Udviklingshistorie, man undersøger de enkelte Arters
Overensstemmelse og grupperer dem derefter; men Spørgsmaalet om selve
Artens Oprindelse drages ikke med ind i Undersøgelsen, fordi den
umiddelbare Iagttagelse ikke viser nogen Overgang fra den ene Art til den
anden.

Dog er der mange Ting, som kunne tyde paa, at Artsbegrebet ikke har
absolut Gyldighed. Ikke blot viser det sig vanskeligt, ja umuligt i
mange Tilfælde at afgjøre, hvad der skal kaldes Art, og hvad der maa
betragtes som blot Varietet; men selve denne Fremstilling om absolut
færdige Livsformer, i Følge hvilken Livet kun var disse Livsformers
Selvopholdelse eller deres Kamp for at hævde sig overfor de ydre Vilkaar
--- denne Fremstilling, hvorved der skarpt sondredes mellem Organismen og
dens Omgivelser, som om de ikke hørte nødvendig sammen som Led i
Livsfænomenernes Kreds, var dog altfor modsigende og udvortes mekanisk,
til at ikke mange Forskere maatte føle Trang til at gaa ud over den.
Saa længe de theologiske Forestillinger endnu beherskede de fleste
Naturforskeres Bevidsthed, kunde denne Trang vistnok tilfredsstilles ved
at føre de absolute Artsforskjelligheder tilbage til forskjellige
Skabelsesakter, idet hver Art tænktes skabt for sig. Men naar
Bevidstheden om naturligt Sammenhæng og naturlige Aarsagers Nødvendighed
udvikler sig, kan Trangen ikke længer tilfredsstilles paa denne Maade.
For saa vidt Naturforskere endnu betragte de forskjellige Arter som skabte
hver for sig, er det i Reglen kun at opfatte som rhetoriske Udtryk for,
at de antage Arternes Oprindelse for uforklarlig. Allerede ved Slutningen
af forrige Aarhundrede fremtræder Ideen om Arternes naturlige Oprindelse,
og mærkværdig nok omtrent paa een Gang i England, Frankrig og Tydskland:
hos Erasmus Darwin (Charles Darwins Bedstefader), hos Geoffroy
St.\ Hilaire og Lamarck og hos Goethe og Oken. Disse ældste Forsøg
henvise vel, hvert for sig, til de ydre Aarsagers Virken, men de fleste
af dem lægge dog Hovedvægten paa en indre Dannelseskraft, en spontan Evne
til Fuldkommengjørelse, som skulde føre de levende Væsener fra lavere til
højere Trin. For saa vidt er der endnu noget Mystisk og Ubestemt i disse
Forsøg.

Det Geniale i den Maade, hvorpaa Darwin har optaget det gamle Problem,
beroer nu paa, at dette Ubestemte og Mystiske træder til Side for
Henviisningen til en bestemt, faktisk virkende Kraft. For saa vidt kan
Newtons Fordring om en \textit{vera causa} siges at være fyldestgjort.
Den Faktor, Darwin henviser til, er det saakaldte Kvalitetsvalg eller
den naturlige Udvælgelse, o: den Kjendsgjerning, at kun de heldig stillede
Organismer leve og forplante sig, medens de svage og uheldig stillede gaa
til Grunde. Naturen renser altsaa de levende Væsener ligesom gjennem et
Sold, og af denne Renselsesproces maa der nødvendigvis fremgaa mere
udviklede og til Betingelserne bedre svarende Former. Det er denne
naturlige Udvælgelse, der frembringer Racer og Varieteter: hvorfor skulde
den saa ikke ogsaa kunne frembringe Arter? »Man har ofte forsikret,« siger
Darwin, »men man kan ikke levere noget Bevis for denne Forsikring, at det
Beløb af Variering, der findes ude i Naturen, har bestemte Grændser
\ldots\ Jeg kan ikke se, at der er nogen Grændse for denne Magt, der
langsomt og skjønt lemper hver enkelt Form efter de mest indviklede
Livsforhold.« I disse Ord finder Darwinismens logiske Gyldighed sit
klareste Udtryk.

Afgjørelsen af Striden beroer nemlig paa det Spørgsmaal: har den
naturlige Udvælgelse Grændser eller ikke? Men dette Spørgsmaal kan
Videnskaben paa sit nuværende Stadium ikke besvare. Det er dette, der
gjør Darwinismen til en Hypothese og retfærdiggjør de ædruelige Forskeres
Betænkelighed ved at se Forsøg paa at udgive den for en godtgjort Lære,
paa hvilken man nu skulde kunne gaa ind. Heraf forklares den franske
Positivismes Stilling til dette Spørgsmaal. Der er uden Tvivl Mange, som
sammenblande Positivisme og Darwinisme og ikke vide, at de udelukke
hinanden. Auguste Comte kritiserede skarpt Lamarcks Theori og saae i den
et uberettiget Forsøg paa at overføre Udviklingens Lov, der gjælder for
den individuelle Organisme, paa Slægten som Helhed. En af hans Disciple,
den berømte Anatom Charles Robin, har nylig gjentaget denne Indsigelse og
viist, at den ogsaa rammer Darwin. Særlig fastholder han, støttet til
Anatomiens exakte Resultater, at Varieringen har sine bestemte Grændser.
Enhver Organisme kan svinge mellem sine Yderpunkter; gaaer Forandringen
ud over dem, udvikles Organismen ikke, men døer. Og netop de levende
Væsener, hvis Organisation er meget simpel, som Zoofyter og Orme, variere
kun indenfor snevre Grændser; men af dem skulde de højere Arter jo have
udviklet sig. Og i ethvert Tilfælde kan man ikke forfølge den ene Arts
Udvikling af den anden, saaledes som man kan forfølge, hvorledes det ene
Stadium af et enkelt Individs Liv udvikler sig af et andet. »Det er
ogsaa Fremskridt og Bevis paa Indsigt at forstaa at undersøge en
almindelig Sætning, førend man antager den.« Robin betragter det derfor
som et Vidnesbyrd om de franske Forskeres filosofiske Aand, at de fleste
af dem erklære sig imod Darwin. Tydskeren Häckel vil rigtignok gjøre det
til en Prøve paa Menneskenes aandelige Overlegenhed, om de antage
Darwinismen og »den derpaa grundede Filosofi« eller ikke, og tillægger
derfor Englændere og Tydskere Æren af at være videre fremskredne i
Kulturudvikling. Med Rette har Littré stemplet dette som Chauvinisme.
Og dersom --- hvad der fra den strenge Videnskabs Standpunkt maa
betragtes som muligt --- dersom den darwinske Hypothese er urigtig, saa
vilde det se betænkelig ud med denne Overlegenhed. Hvad Kant tvivlede
om --- at der nogensinde vilde opstaa en Newton for Videnskaben om det
organiske Liv --- det seer Häckel virkeliggjort i Darwin. Han mener, at
Darwins Theori er en af Menneskeaandens største Erobringer, der
»umiddelbart kan stilles ved Siden af Newtons Gravitationslære« --- ja
at den endog staaer over denne! Imod en saadan Overseen af den induktive
Methodes Principer, som her lægges til Grund, maa Logiken protestere.
Newtons Opdagelse staaer som et Ideal i Videnskaben, fordi den er et
enestaaende Exempel paa Forening af genial Skuen og streng Bevisførelse,
hvor intet Led er oversprunget, men alle Enkeltheder føje sig sammen til
en storartet Bygning. Darwin har paavist en real Aarsag, som er af
væsenlig og maaske hidtil overseet Betydning for det organiske Liv, men
han har ikke kunnet bevise, at denne Aarsag indeholder den tilstrækkelige
Grund til de forskjellige Arters Oprindelse.

Vil det nu være vanskeligt og langvarigt at finde et afgjørende Svar
paa det Spørgsmaal, om Varieringen har bestemte Grændser eller ikke?
Darwin har selv betraadt den Vej, som ene kan føre til et exakt Bevis,
nemlig Experimentets Vej. Som bekjendt var det Husdyrene og de dyrkede
Planter, mod hvilke Darwin især vendte sin Opmærksomhed, og hvorfra han
hentede sit vigtigste Materiale. Dersom man nu kunde frembringe
tilstrækkelig varierende Betingelser, og i tilstrækkelig lang Tid, saa
vilde Spørgsmaalet maaske finde sin Besvarelse; men det er jo slet ikke
sagt, at det staaer i Menneskets Magt at tilvejebringe de Betingelser,
hvorved i den frie Natur Arter maatte have udviklet sig af hverandre.
Dette er altsaa kun en meget ubestemt og fjern Mulighed.

En væsenlig, i Følge Darwins egen Mening den væsenligste Indvending mod
hans Hypothese hentes fra, at »endskjøndt de geologiske Undersøgelser
utvivlsomt have lært os mange tidligere Leds Tilværelse at kjende, Led,
som bringe talrige Livsformer meget nærmere til hinanden, saa have de
dog ikke givet de uendelig mange fine Gradationer, som Theorien fordrer,
mellem uddøde og nulevende Arter.« Ogsaa her har det lange Udsigter med
et exakt Bevis, ja det er maaske tvivlsomt, om et saadant nogensinde kan
naaes ad denne Vej, eftersom ethvert nyt Mellemled, der opdages, stedse
vil kunne betragtes som en ny Art, der kan stikkes ind mellem to andre;
selve Overgangen fra een Art til en anden kan man efter Sagens Natur ikke
finde i de geologiske Lag, der kun indeholde de døde Rester.

Ganske vist vil Sandsynligheden af Hypothesens Sandhed stige, jo flere
saadanne Mellemled der paavises. Man seer ogsaa, at Forskere, som
tidligere vare Modstandere af Arternes Foranderlighed, opgive deres
Modstand under Indtrykket af, hvad de nye Fund bringe for Dagen. »At en
Naturforsker, der saa længe og saa grundig har studeret de talrige
Levninger og sammenlignet dem saa omhyggelig og utrættelig baade med
beslægtede Levninger fra andre Steder og med de nulevende Arters Skeletter
som Gaudry, har tabt Troen paa Arternes absolute Fasthed og slutter sig
til dem, der tro paa deres Omdannelse i Tiden, er under alle
Omstændigheder et betydningsfuldt Fingerpeg, saa meget mere, som han
neppe kan have været paavirket i sin Dom af de Naturforskere, i hvis
Nærhed han arbejdede; thi ingensteds har Darwinismen vundet mindre Indgang
end i Frankrig.«

Men et exakt Bevis synes der, som sagt, ikke at være nogen Udsigt til, i
det Mindste i lange Tider. For saa vidt have de, der af moralske og
religiøse Grunde tro at maatte kæmpe mod Darwinismen, Ret, naar de
spottende henvise til det Tidspunkt, »da Darwinismen bliver bevist.« Dog
turde deres Glæde alligevel være noget forhastet. Eet er det nemlig, at
de Faktorer, hvoraf man hidtil har villet forklare Arternes Oprindelse,
ikke strække til, et ganske Andet er det, at Arternes Oprindelse i det
Hele betragtes som et Spørgsmaal, der hører hjemme i Naturvidenskaben og
ikke i Theologien. Og dette Sidste synes i ethvert Tilfælde at være
opnaaet ved den Darwinske Strid. At Arterne maa have en naturlig
Oprindelse, at Aarsagerne til deres Opstaaen, hvilke de saa ellers ere,
maa søges i de i Naturen givne Betingelser, det er en Forudsætning, som
udspringer af Principerne for al videnskabelig Forsken, Principer, som
mere og mere komme til Anvendelse selv paa Omraader hvorfra man tidligere
holdt dem ude. Den Grundsætning, at enhver Ting har sin Aarsag,
udspringer af selve den menneskelige Bevidstheds inderste Lov og kan ikke
trænges tilbage, naar den først er bleven anerkjendt. Hvad Aarsagen er,
det er et andet Spørgsmaal. Mange Hypotheser af lignende Art som Darwins
ville maaske afløse hverandre, førend den sande findes. »Det er,« siger
Mill, »ingen gyldig Grund for Antagelsen af en given Hypothese, at vi
ere ude af Stand til at forestille os nogen anden, der kan forklare
Kjendsgjerningerne. Der er ingen Nødvendighed for at forudsætte, at den
sande Forklaring maa være en, som vi med vor nuværende Erfaring kunne
forestille os.« Men i Principet staaer det dog fast, at der maa være en
naturlig Aarsag. Den franske Positivismes Skepsis er kun begrundet
overfor de enkelte Forklaringsforsøg, ikke hvad selve det Principielle
angaar. Det kan være forhastet allerede nu at ville løse dette Spørgsmaal,
men det hører ikke til dem, som i sig selv ligge ud over Videnskabens
Grændser.

\medskip
\begin{center}
[\emph{Sluttes.}]
\end{center}

% ============================================================
%  ANDEN DEL (Nær og Fjern, nr. 94)
% ============================================================

\bigskip
\begin{center}
\textbf{Anden Del}
\end{center}
\medskip

\section*{II.}

Den logiske Undersøgelse af Darwinismen fører saaledes til at anvise den
dens rette Plads som en genial Hypothese, der er fremgaaet af den nyere
Tænknings og Videnskabs Princip om Naturlovenes og de naturlige Aarsagers
Herredømme i Tilværelsen. Dette Princip har ikke ventet paa Darwinismen
for at træde i Kraft; det har fundet sin logiske og metafysiske Begrundelse
i den nyere Filosofi og sin reale Gjennemførelse i den nyere
Naturvidenskab. Det staaer og falder derfor heller ikke med Darwinismen.
Vi kunne altsaa ikke tillægge denne en saa stor Betydning for hele
Verdensanskuelsen, som man ofte gjør; den er et værdifuldt, om end
hypothetisk Led i en større Sammenhæng.

Dog kan den siges at være bleven ligesom et Stikord i Kampen mellem den
supranaturalistiske og spiritualistiske Opfattelse og den humane og
rationelle Verdensanskuelse. De Indvendinger, der rettes imod Darwinismen
ud fra hin Opfattelse, vilde blive rettede mod en hvilken som helst anden
Lære, der vilde forklare Arternes Oprindelse paa naturlig Maade. Derfor
vil det, aldeles afset fra Darwinismens hypothetiske Karakter, have
Interesse at undersøge, i hvilken Retning den peger, naar man drager dens
sidste Konsekvenser.

Den almindelige Betydning lyder paa, at Darwinismen er en materialistisk
Lære, som forklarer det Indre af det Ydre, det Aandelige af det Legemlige,
og opfatter Mennesket som et Dyr, idet det lader ham »nedstamme fra
Aberne«. Istedet for den guddommelige Visdom, der efter den gamle Tro
leder Alt i Naturen og i Historien til det Bedste, seer man at den blinde
Tilfældighed raader, det tilsyneladende Hensigtsmæssige skyldes kun en
lykkelig Hændelse eller rettere en Række af lykkelige Hændelser.

Vi tage her ikke Hensyn til Darwins personlige Mening, uagtet den ikke
synes at gaa i materialistisk Retning. For at prøve den anførte Betydning
holde vi os til den Aarsag, hvoraf han vil forklare Arternes Oprindelse,
nemlig Kampen for Tilværelsen. Det er for det Første klart, at Kamp
forudsætter en Kæmpende og Vaaben og Kraft hos denne Kæmpende. Det vil
sige, der forudsættes en indre Activitet i Organismen og et oprindeligt
Anlæg, hvorfra Udviklingen kan tage sin Begyndelse. Uden dette hjælpe de
ydre Paavirkninger Intet. Dernæst ligger der i Udtrykket: Kamp for
Tilværelsen ogsaa det, at Tilværelsen er et Gode, som Individet stræber
at bevare. Thi det er aabenbart ikke blot det at være til, det er den med
Tilværelsen forbundne Nydelse og det Værd, Tilværelsen derved faaer,
hvorom der kæmpes. Men disse Forudsætninger vise ikke i materialistisk
Retning. De føre tvertimod til Anerkjendelsen af ideelle Principer som de
egentlig afgjørende. Darwin selv har dog vistnok lagt for stor Vægt paa
de ydre Aarsager. Den berømte Botaniker Nägeli har i Modsætning hertil
udtalt som Resultatet af sine Undersøgelser, at Dannelsen af Afarter og
Racer ikke er en Følge af og et Udtryk for ydre Paavirkninger, men
betinges ved indre Aarsager, og han tilskriver derfor Organismen et
»Fuldkommengjørelsesprincip«, som stadig fører til nye Former, en Proces,
under hvilken de ydre Vilkaar kun spille Regulatorens Rolle. Vi føres
herved tilbage til Livets almindelige Lov, saaledes som af de nyere
Filosofer Herbert Spencer skarpest har formuleret den: en Tillempelse af
det Indre efter det Ydre. Det Indre er altsaa stedse forudsat som en
oprindelig Faktor, der gjennem sin Vexelvirkning med Omverdenen arbejder
sig fremad til fuldkomnere Former. I Følge Spencer har Darwin lagt for
ringe Vægt paa den direkte Tillempelse, som skeer ved Funktionens
Indflydelse paa Strukturen. Kun naar Organismen ikke formaaer at lempe
sig efter Omgivelserne og derved overvinde dem, træder den indirekte
Tillempelse i Kraft, som skeer ved de Svages Udryddelse.

Hvilken er altsaa den Verdensanskuelse, man maatte uddrage af Darwinismen,
dersom den viste sig at have Ret? Ingen anden end den, der fremgaaer af
den nyere Tids hele filosofiske og videnskabelige Bevidsthed, som ikke
kan antage Ideelle Magter, der skulde virke paa anden Maade end gjennem
selve de reale Betingelser og de Love, der gjælde for disse. Vi kunne,
som en aandrig Tænker har sagt, tro paa virkende, men ikke paa herskende
Ideer. Det Ideale er i sig selv ikke fremmed for de reale Betingelser;
disse ere Elementer i dets Væsen, og det godtgjør netop sin Guddommelighed
ved at arbeide sig frem gjennem alle Hindringer. Den Ide, vi tro paa, maa
være en kjæmpende Ide.

Det er saa langt fra, at Darwinismen fører til Materialisme, at den
netop maa siges at vise i den helt modsatte Retning. Den lægger ganske
vist en stærk Vægt paa de ydre Vilkaar; men disse Vilkaar virke kun
gjennem Organismens Lempelse efter dem. Og hvad man kalder Tilfælde, er
kun Exemple paa Betingelsernes, de indre og ydre Betingelsers, uendelig
forgrenede og indviklede Sammenhæng, som vi ikke kunne overse; i de
saakaldte tilfældig opstaaede Former have vi netop et Bevis paa Livets
Energi, der veed at tilegne sig endog de fineste Nuancer i de ydre
Vilkaar og anvende dem i sin Tjeneste.

Men, siger man, Darwins Lære er trøstesløs, fordi den viser os, at de
Svage gaa til Grunde og kun de Stærke og Lykkelige bestaa. Hvor
skrækkeligt at vide sig og Alt prisgivet en saadan mørk og ubønhørlig
Magt! Darwin selv har følt al den Lidelse og Smerte, som
Tilværelseskampen fører med sig. »Vi se Naturens Aasyn straalende af
Glæde, vi se ofte, at der er Næring i Overflod, men vi se ikke, eller vi
forglemme, at de Fugle, der sidde og kvidre rundt om os, for Størstedelen
leve af Insekter eller Frø og saaledes til enhver Tid ødelægge Liv; og
heller ikke mindes vi, i hvor høj Grad disse Sangere, deres Æg eller
deres Unger ere udsatte for Ødelæggelse af Rovfugle og Rovpattedyr; vi
have heller ikke altid det staaende klart for os, at skjøndt der i
Øjeblikket kan være mere end Føde nok, er det dog ikke saaledes hvert
Aar eller til hvilken som helst Aarstid.« Men han trøster sig med, at
»af Naturens Kamp, med Hunger og Død, fremgaaer det Høieste af det, vi
kunne opfatte.«

Det er ikke let at indse, hvilken Fordel den almindelige Tro paa
»Hensigtsmæssigheden i Naturen« her har for Darwins Lære. Thi hvilken
skarp Strid er der dog ikke mellem den almægtige og forudseende Visdom,
som skal have frembragt Alt, og den Smerte og Nød, den Lidelse, alt
Levende er underlagt, og som hverken Darwinister eller Antidarwinister
kunne overse! Her er Noget, som ingen Theori, ingen Verdensanskuelse kan
forklare. Det Høieste, vi kunne naa, er at fæste vort Blik paa Maalet,
der naaes gjennem disse Kampe, og se, hvorledes at hin Smerte er en
Følge af Udviklingen hen imod dette Maal. Ved Siden af den Forvisning,
at denne Udvikling ustandset gaaer sin Gang, viser Resignationens
Nødvendighed sig; vi se de bestemte, ufravigelige Betingelser, uden hvilke
Intet kan naaes her i Verden, og som vore Ønsker og Længsler nu en Gang
ikke kunne forandre. Men denne Resignation maa en hvilken som helst
Verdensanskuelse nødvendigvis under en eller anden Form forlange.

Ligesom Kants og Laplaces astronomiske Theori og Lyells geologiske Theori,
saaledes lærer Darwins biologiske Theori os, at de samme Kræfter, som i
ethvert Øieblik umærkeligt virke rundt om os, ved deres stadige, aldrig
afbrudte Virken have frembragt de Former i Naturen, vi i den Grad
beundre, at vi ere tilbøjelige til at tænke os dem frembragte i et Nu,
ved en almægtig Skabelsesakt. Ligesom Kopernikus aabnede Universets
Uendelighed for vort Blik, gjorde Ende paa den gamle Modsætning mellem
Himmel og Jord, og førte os til at betragte vor Klode som et forsvindende
Punkt i Verdeners Uendelighed, saaledes føre hine nyere Theorier os ind i
Kræfternes Uendelighed og Evighed. De mest afgjørende Modsætninger ere
kun relative Stadier i den uendelige Proces. Vi kunne ikke omfatte det
Uendelige; men alle de Forskjelle, vor Tanke vil statuere, løse sig
stedse op paa Ny, vise sig at være relative. I Kræftens og Rummets
Uendelighed har Aandens Uendelighed sit Symbol. Den højeste Ide, vi
kunne danne, er Tanken om en uendelig Verdensharmoni, hvor ethvert Enkelt
paa sit Sted og til sin Tid spiller sin Rolle, gjennem sin Kamp og
Lidelse yder et realt Bidrag til den store Livsproces.

Naar man seer, hvilken Begejstring det kopernikanske System vakte, og
hvorledes det gav Anledning til den første Anvendelse af
Naturvidenskabens Resultater paa Livsanskuelsen, saa kan man ogsaa vente
lignende Følger af disse nyere Theorier, saa sandt som Kræfternes
Uendelighed ikke er mindre opløftende end Verdensrummets Uendelighed.
De overleverede Troeslærdomme modsætte sig Kopernikus og Newton, som de
nu modsætte sig Lyell og Darwin; men som de have maattet give Kjøb
overfor hine, saaledes ville de ogsaa komme til at gaa paa Akkord med
disse. Vi ere nu en Gang nødte til at danne vore højeste Ideer efter den
Maade, hvorpaa Universet fremstiller sig for os. Vor Livsanskuelse er en
Verdensanskuelse. Ligesom Kopernikus' og Newtons Opdagelser nøde Tanken
til at lade sin højeste Ide omfatte Verdensuendeligheden, hvilket allerede
Giordano Bruno indaandede med en Overbevisning og Begejstring, der end
ikke svigtede paa Inkvisitionens Dødsbaal, --- saaledes maa nu i Følge
Udviklingteorien Opfattelsen af den højeste Ides Virkemaade undergaa en
tilsvarende Forandring. Ikke ved pludselig Indgriben, ikke gjennem
uforberedte Katastrofer og Revolutioner, men stille og umærkt virker
den evige Kraft, i hvis Væsen alle Betingelser, alle reale Aarsager ere
Elementer.

Jeg har hørt en begavet Dame sige, at dersom Darwin havde Ret, og vi
virkelig nedstammede fra Aberne, vilde Livet tabe alt Værd og Betydning
for hende. Det er maaske dette Punkt, som for de Fleste staaer som det
mest afskrækkende. Men den, som seer, at selv det Højeste og Ædleste dog
har sine bestemte Betingelser, at det ikke kan være uden dem, kun er ved
at tilegne sig dem, --- den, som erkjender, at de tilsyneladende absolute
Forskjelle forsvinde i den uendelige Sammenhæng, vil se Sagen paa en
anden Maade. Han vil ogsaa i Menneskeaanden se en Form for den kæmpende
Ide og vil ikke se nogen Nødvendighed for, at det Aandelige ude- eller
ovenfra indblæses i Materien; det Aandelige staaer ikke som et fremmed
Element i Tilværelsen; det fremkommer paa sit bestemte Trin i
Udviklingen, arbejder sig frem af selve Materien og afgiver i sidste
Instans den eneste Forklaring af denne. Det Paafaldende i den
Fremstilling, at Menneskene skulde nedstamme fra Aberne, har sin simple
Grund deri, at man her umiddelbart stiller to Udviklingstrin sammen, der
ligge uendelig langt fra hinanden; man overspringer da netop det
Væsenlige --- selve Udviklingen. Hvor stor en Kløft er der ikke mellem
de Menneskeædere, fra hvilke vi maaske nedstamme, og vor Tids
civiliserede Mennesker!

Vi se det enkelte Menneskes Aand udvikle sig fra en ringe Begyndelse; vi
se i Historien Menneskeslægten udvikle sig fra dyrisk Raahed til aandelig
Dannelse og humant Samfundsliv. Hvad Darwinismen forlanger, er altsaa
kun, at man skal tænke sig de Linier, denne Udviklingsgang beskriver,
forlængede i Retningen nedad, ligesom vi i alle vore ideale Drømme, i
vor Stræben efter Fremskridt og højere Udvikling tænke os dem forlængede
i Retningen opad. Dette fører os til det sidste Punkt, vi ville fremdrage
til Belysning af Darwinismen, dens ethiske Karakter. Men førend vi
forlade den almindelige Verdensanskuelses Omraade, vil jeg som et
\textit{argumentum ad hominem} minde om, at en dansk Filosof, som for os
Alle staaer som en sjelden ædel og human Tænker, og hvis Livsanskuelse
havde en særegen Varme og Inderlighed --- at afdøde F.\ C.\ Sibbern
nærede den samme Anskuelse om Menneskets Oprindelse, som man nu forarges
saa højlig over, naar man møder den hos darwinistiske Forfattere. Det
hedder nemlig i hans Skrift »Spekulativ Kosmologi«

(Kjøbenhavn 1846) P. 20--21: »Er det første Menneske
blevet til paa anden Maade end vi, saa er det dog ej
skabt i nogen anden Forstand af Ordet »skabt« end vi.
Hvorledes mon det første Menneske blev til? Uden Tvivl
derved, at en vis Dyreformation fra en første ringe
Begyndelse, gjennem en Række af Generationer eller
Afødninger, idet i Afkommet bestandigen Uddannelsesprocessen
gik videre og videre, naaede til det Punkt, da nu dets
sidste Affødning blev et Væsen, som saa at sige slog
Aandelivhedens Øje op, idet det sagde Jeg til sig selv
og følte sig i sin Jeghed, men nu og følte Drift i sig
til at navngive alle Dyr, d.\,e.\ til at stemple sig dem,
samt til at betragte sig som staaende i Relation til dem
alle og endog over dem alle. Hermed gik Menneskeheden
op af Dyriskheden.«

\section*{III.}

Ligesom man i Almindelighed mener, at Darwinismens
Metafysik er Materialisme, saaledes mener man, at dens
Ethik maa gaa ud paa dyrisk Sandselighed. »Lader os æde
og drikke, thi i Morgen skulle vi dø!« --- det er den
Tankegang, man i Reglen vil anse for den naturlige hos en
Darvinist. Mange gaa endnu videre og mene, at Selvmord
er den naturlige Ende for den, der nærer saadanne
Anskuelser. »Ethvert Selvmord,« skrev en amerikansk Præst
til Tyndall, »ethvert Selvmord i vort Land (og de
indtræffe daglig) er en indirekte Virkning af de bestialske
Lærdomme, som De, Darwin, Spencer, Huxley et id omne
genus forkynde.«\footnote{Brevet er aftrykt i H.\ Spencers:
\textit{The Study of Sociology}, pag.\ 419.}

Der er neppe noget større Bevis paa den snævre
aandelige Horisont, de fleste Mennesker endnu besidde, end
den store Vanskelighed, de have ved at tænke sig, at der
kan leves et moralsk og ædelt menneskeligt Liv paa Basis
af andre Livsanskuelser end deres egen. Den anderledes
Troende er for de Fleste et umoralsk Væsen. Og bag
Beskyldningen for Umoralitet ligger Noget, som er værre;
der ligger Fanatismens Had. Saaledes skriver hin
amerikanske Præst fremdeles i sit Brev: »Skulde jeg ikke
hade dem, som hade Dig, Herre!«

I Modsætning til de Lidenskaber, Kampen sætter i
Bevægelse, maa vi holde os til Sagen selv og undersøge
den. Darwinismen er kun en særegen Form eller Anvendelse
af Udviklingens almindelige Lov, der før Darwins Tid var
bleven opdaget som gjældende ikke blot paa Naturens
Omraade (Kant, Laplace, Lyell), men ogsaa paa Historiens
Omraade. Allerede hos Baco af Verulam finde vi
Forestillingen om Menneskehedens Udvikling og
Fremadskriden. Pascal og Leibnitz udtrykte denne
Forestilling nøjere. Men det er i forrige Aarhundrede, at
de første Skridt gjøres til dens Gjennemførelse af Turgot
og Condorcet, Lessing og Kant. Ethiken, der tidligere kun
havde fæstet sit Blik paa den enkelte Personligheds
Udvikling, maatte nu faae Øje for, at denne Udvikling kun
kan foregaa som Led i Slægtens almindelige Udvikling.
Udvikling af det personlige Liv gjennem Deltagelse i det
kulturhistoriske Arbejde, det er det Moralprincip, som har
sin Kilde i Udviklingens Ide.

Darwinismen har de samme ethiske Motiver som hele den
paa denne Ide byggede Verdensanskuelse. Selve den store
Rolle, Udviklingen spiller i Darwins Hypothese, viser
i ethisk Retning. Hvile, Stillestaaen er ikke ethisk; det er
Arbejde, Vorden, Kamp og Fremgang, det gjælder om, og
dertil fører Darwinismen os umiddelbart, idet den indprenter,
at den, der ikke kæmper af al sin Magt, vil bukke under,
og at Livet selv er en Kamp for Tilværelsen. Umiddelbart
seet er det vistnok den blinde, egoistiske Selvopholdelsesdrift,
der kommer til at spille den største Rolle; men den vil
efterhaanden lige som aabne sig, jo mere Udviklingens Gang
fører Individerne sammen og lærer dem at opfatte den
fælles høiere Lov, for hvilken de enkelte Villier maa bøje
sig. Selvopholdelsesdriften lutres da til en ideal Stræben
efter at bevare og udvikle den sande menneskelige Tilværelse
baade hos Andre og hos sig selv. Saaledes bliver den
højeste ethiske Stræben endnu en Kamp for Tilværelsen, og
den Naturlov, Darwin har paavist, faaer sin Gyldighed
ogsaa paa det ethiske Omraade.

Dersom Darwin havde Ret i sin Antagelse, at Varieringen
ingen Grændser har, vilde det endogsaa være en storartet
og begejstrende Udsigt, han aabnede for al højere Stræben.
Mennesket vilde da ikke støde paa Grændser for sit ideale
Higen; hvad han ikke selv kunde naa, vilde han kunne
lægge Grunden til for Slægten efter ham. Og ethvert nok
saa lille Fremskridt, enhver nok saa uanseelig og
tilsyneladende ubetydelig Virken vilde dog give sit Bidrag
til Udviklingen; det er jo gjennem en uendelig Mængde
uendelig smaa Virkninger, at de store Resultater opnaaes.

Paa den anden Side indeholder Darwinismen en Lære,
som den tidligere spiritualistiske Ethik ofte har overseet,
nemlig om de ydre Vilkaars store Betydning. Arbejdet for
Menneskislægtens Fremgang bestaaer ikke mindst i Forbedring
af dens materielle Forhold. Darwinismen viser os, hvilken
Betydning disse Forhold have for en Arts Udvikling, altsaa
ogsaa for Menneskehedens. Det bliver altsaa en ethisk
Opgave at skaffe de fornødne Udviklingsbe\-tingelser tilstede
for Alle. Darwinismens ethiske Konsekvenser pege saaledes i
Retning af den sociale Reform. Den viser os, hvilken Vildhed
og hvilken Rædsel der opstaaer, naar Kampen drejer sig om
den materielle Selvopholdelse, og indskærper Nødvendigheden
af at bringe det til, at den Kamp, uden hvilken
Menneskelivet vilde stagnere og udslukkes, mere og mere
kan føres paa en Menneskeheden værdig Maade, og om
Formaal, der hæve sig over de blinde dyriske Drifter, der
kun virke, fordi Nøden fremkalder dem, og fordi de ikke
trænges tilbage af højere Ideer.

Vort Resultat er altsaa følgende. Darwinismen er en
Hypothese, og det er ulogisk og uvidenskabeligt at betragte
den som Mere. Men den appellerer ikke blot til en virkelig
og betydningsfuld Aarsag; den er selv fremgaaet af den
nødvendige Forudsætning, at det ogsaa ved Arternes
Oprindelse maa have været naturlige Aarsager, der virkede.
De mislige theoretiske og praktiske Konsekvenser, man har
antaget for en nødvendig Følge af den, ere urigtige, og den
Modstand, der gjøres mod Darwinismen, gjælder egentlig
hele den rationelle og humane Verdensanskuelse, som ikke
staaer og falder med Darwinismen. Denne Verdensanskuelse
var i sine afgjørende Grundtræk udviklet, før Darwinismen
optraadte, og vil vedblive at udvikle sig, selv om
Darwinismen skulde blive gjendreven. Men paa den anden
Side vilde dennes Sandhed skaffe hin Verdensanskuelse en
betydningsfuld Tilvæxt, berige den med nye
Synspunkter og Ideer. --- Men, som sagt, den som veed,
haster ikke!

\end{document}
